\chapter{Electron Trajectories and Categorical Measurement}
\label{ch:electron_trajectories}

\papernote{Source paper: \emph{On the Consequences of Hardware Oscillatory Dynamics: Multi-Modal Ensemble Virtual Spectrometry Methods for Electron Trajectory Completion} --- \texttt{docs/electron-trajectories/}}

\begin{chapterabstract}
Electron trajectories during atomic transitions are observable without wavefunction collapse through categorical measurement of partition coordinates. The commutation relations establish that categorical observables are orthogonal to physical observables by construction: measuring categorical position does not disturb momentum. The quintupartite ion observatory---five simultaneous spectroscopic modalities---achieves trans-Planckian temporal resolution at $10^{-138}$~seconds, 95~orders of magnitude below the Planck time. Classical mechanics, thermodynamics, and electromagnetism are all derived as limits of the categorical framework: they are observational perspectives on partition geometry, not independent theories.
\end{chapterabstract}

\input{../docs/electron-trajectories/electron-trajectories.tex}
